\section{METHODS}\label{methods}

This section defines the canonical loss-grid algorithm and the
three optimization approaches mapped to the research questions:
algorithm-level optimization candidates
(Section~\ref{methods-redesign}, RQ1), multi-device
heterogeneous scheduling (Section~\ref{methods-scheduling},
RQ2), and adaptive calibration with cache reuse
(Section~\ref{methods-calibration}, RQ3). It also states the
shared benchmarking, surface-validation, and uncertainty
protocol. Terminology is defined in
Section~\ref{introduction}; experimental procedures in
Section~\ref{experiments}.

\subsection{Canonical Loss-Grid Algorithm
(Baseline)}\label{methods-canonical-algorithm}

The thesis takes Li et al., \emph{Visualizing the Loss
Landscape of Neural Nets}~\cite{li2018losslandscape}, as the
canonical method. It is the de-facto reference in subsequent
work~\cite{xie2025losslens,losslens_repo} and the algorithm
shipped to end-users~\cite{goldstein_loss_landscape_repo};
alternative formulations (e.g.\ trajectory sampling) probe a
different aspect and are not the cost driver in
LossLens-style workflows.

For one grid point the algorithm picks coordinates
$(\alpha,\beta)$, perturbs parameters via $w(\alpha,\beta) =
w_0 + \alpha\,d_1 + \beta\,d_2$, evaluates loss over the
dataset subset, and writes the scalar into the surface; this
repeats until the grid is complete. Here $w_0$ is the trained
parameter vector and $d_1, d_2$ are filter-normalized
directions following the upstream default \texttt{weights}
mode~\cite{li2018losslandscape,goldstein_loss_landscape_repo}.

The thesis adopts the procedure as a controlled single-machine
comparison baseline with two adjustments to align with the
research questions. First, the baseline uses one GPU, matching
the problem framing in Section~\ref{problem-statement} and
leaving RQ2 to measure CPU--GPU scheduling rather than
GPU--GPU distribution. Second, the baseline performs no CPU
helper execution, because hybrid execution is itself the RQ2
candidate; including it in the baseline would make that
comparison circular. Because the baseline is single-GPU, the
implementation runs in a single process rather than through
the public MPI/distributed path. The only omission relative
to the public MPI
reference~\cite{goldstein_loss_landscape_repo} is distributed
execution. The per-grid-point semantics are preserved, but
grid points are evaluated sequentially. All comparisons in
Section~\ref{results} are relative to this baseline. A
perturbed setting $w(\alpha,\beta)$ is a \emph{perturbation}
of a checkpoint; \emph{checkpoint} is reserved for a distinct
trained model from a separate run.

\subsection{Algorithm-Level Optimization Candidates
(RQ1)}\label{methods-redesign}

The canonical loop evaluates one perturbed parameter state per
grid point against the same data subset. This is the same
outer-loop shape that PyTorch's model-ensembling
tutorial~\cite{pytorch_ensembling} addresses for stacks of
independently trained models, via the composition of
\texttt{functional\_call}~\cite{pytorch_functional_call} and
\texttt{vmap}~\cite{pytorch_vmap}.
\texttt{torch.compile}~\cite{pytorch_compile} is an orthogonal
intra-device mechanism — ahead-of-time graph capture, kernel
fusion, and dispatch elision — on the same per-grid-point
work. RQ1 tests three candidates against the baseline
(Table~\ref{tab:opt-map}): \texttt{vmap} composed with
\texttt{functional\_call} (the documented ensembling pattern),
\texttt{torch.compile} on the baseline evaluation path, and
the composition of both.
\texttt{functional\_call} is a prerequisite primitive rather
than a standalone candidate: it makes no independent
optimization claim, and the candidates that build on it
subsume its observable runtime contribution.

The per-candidate verdict is the effect-size CI on the
$T_{\text{grid}}$ speedup ratio
(Section~\ref{methods-benchmarking}). It is paired with an
\emph{intra-axis composition verdict}
$\in \{\text{multiplicative}, \text{redundant},
\text{destructive}, \text{inconclusive}\}$, derived by
comparing \texttt{compiled\_vmapped}'s observed speedup to the
product of \texttt{vmapped}'s and \texttt{compiled}'s
individual speedups under the same effect-size CI rule on the
ratio.

\begin{table}[t]
\caption{Algorithm-level optimization candidates.}
\label{tab:opt-map}
\centering
\small
\begin{tabular}{@{}p{0.27\linewidth}p{0.27\linewidth}p{0.36\linewidth}@{}}
\toprule
Candidate & Components & Method role \\
\midrule
baseline & in-place mutation + sequential loop & reference implementation \\[4pt]
\texttt{vmapped} & \texttt{functional\_call} + \texttt{vmap} + chunking & vectorized forward/loss over perturbations \\[4pt]
\texttt{compiled} & \texttt{torch.compile} on baseline eval path & graph capture and kernel fusion on per-grid-point work \\[4pt]
\texttt{compiled\_vmapped} & \texttt{torch.compile} over the \texttt{vmapped} candidate & intra-axis composition: graph capture over vectorized point evaluation \\
\bottomrule
\end{tabular}
\end{table}

\subsubsection{\texorpdfstring{Vectorized Point Evaluation with
\texttt{torch.vmap}}{Vectorized Point Evaluation with torch.vmap}}\label{vectorized-point-evaluation-with-torch.vmap}

Loss-grid computation has the same outer-loop shape as model
ensembling: a stack of parameter states evaluated against the
same data. The model-ensembling
tutorial~\cite{pytorch_ensembling} composes
\texttt{vmap}~\cite{pytorch_vmap} with
\texttt{functional\_call}~\cite{pytorch_functional_call} to
evaluate that stack in one batched call.
\texttt{functional\_call} replaces in-place parameter mutation
with a stateless invocation: the perturbed parameter mapping
is supplied to the model call for that invocation, leaving
the module object unrewritten. \texttt{vmap} then applies the
loss computation across a leading batch dimension over the
stacked parameter tensors. The \texttt{vmapped} candidate
instantiates this composition: it stacks the perturbed
parameter tensors for $K$ grid points into batched tensors of
shape $(K,S)$ and applies the loss computation across the $K$
slices per data batch
(Figure~\ref{fig:point-evaluation-paths}).

\begin{figure}[t]
\centering
\begin{minipage}{0.96\linewidth}
\small
\textbf{Functional parameter application for grid point $i$}
\begin{algorithmic}[1]
\Statex \textbf{Baseline path}
\State $\mathit{model}.\mathit{parameters} \gets \mathit{params}_i$
\State $\mathit{loss}_i \gets \mathit{model}(\mathit{batch})$
\Statex
\Statex \textbf{Functional path}
\State $f_i \gets \textsc{FunctionalCall}(\mathit{model}, \mathit{params}_i)$
\State $\mathit{loss}_i \gets f_i(\mathit{batch})$
\end{algorithmic}

\vspace{0.6em}

\textbf{Vectorized point evaluation over the grid}
\begin{algorithmic}[1]
\Statex \textbf{Baseline path}
\For{$\mathit{point} \in \mathit{grid}$}
  \For{$\mathit{batch} \in \mathit{data}$}
    \State $\mathit{loss}(\mathit{point}, \mathit{batch})$
  \EndFor
\EndFor
\Statex
\Statex \textbf{Vmapped path}
\For{$\mathit{point\_chunk} \in \mathit{grid\_chunks}$}
  \For{$\mathit{batch} \in \mathit{data}$}
    \State $\mathit{losses}(\mathit{point\_chunk}, \mathit{batch})$
  \EndFor
\EndFor
\end{algorithmic}
\end{minipage}
\caption{Baseline, functional, and vmapped point-evaluation paths.}
\label{fig:point-evaluation-paths}
\end{figure}

The predicted mechanism, if the pattern transfers, is dispatch
amortization and improved use of backend batching rules;
actual kernel behavior is backend- and architecture-dependent.
Full-grid vectorization may exceed memory, so the candidate
evaluates chunks of $K$ grid points at a time and sweeps $K$
for breadth.

\subsubsection{\texorpdfstring{Graph Capture with
\texttt{torch.compile}}{Graph Capture with torch.compile}}\label{graph-capture-with-torch.compile}

\texttt{torch.compile}~\cite{pytorch_compile} captures the
forward/loss path as an ahead-of-time graph, fuses kernels
where the backend permits, and elides Python-level dispatch.
The \texttt{compiled} candidate applies \texttt{torch.compile}
to the baseline per-grid-point evaluation path, isolating the
graph-capture/fusion mechanism from any batching effect; the
loop structure is unchanged. The \texttt{compiled\_vmapped}
candidate applies \texttt{torch.compile} to the
\texttt{vmapped} candidate's batched function, testing
whether graph capture and \texttt{vmap}'s batching multiply on
the same work or interfere. PyTorch documents no canonical
composition of \texttt{torch.compile} with \texttt{vmap} on
the loss-grid pattern, so the composition verdict is
empirical.

\subsection{Multi-Device Heterogeneous Scheduling
(RQ2)}\label{methods-scheduling}

The second approach evaluates heterogeneous CPU/GPU scheduling
for parallel grid-point evaluation. Its premise inherits the
CPU/GPU throughput ratio as the core variable of
heterogeneous CPU/GPU scheduling applicability, identified by
Gallet and Gowanlock~\cite{gallet2021heterogeneous} for set
operations: near or above parity, the CPU can absorb a
non-trivial share of tasks; far below parity, the GPU-only
baseline is already near-optimal. The scheduling mechanism is
dynamic self-scheduling of Polychronopoulos and
Kuck~\cite{polychronopoulos1987gss}, the simpler variant of
the StarPU class of heterocompute
schedulers~\cite{augonnet2011starpu}. The simpler variant is
sufficient because loss-grid is embarrassingly parallel with
uniform per-task work, removing the load-imbalance and
heterogeneous-task-cost problems that motivate the richer
StarPU policies.

The scheduler is eager: GPU and CPU workers poll a shared
queue, one grid point per task, with the same
surface-validation contract as the baseline. Both devices run
the same forward+loss code, keeping the CPU/GPU throughput
ratio interpretable as RQ2's predictor. CPU workers are
bounded by the available logical cores. One grid point per
task is the natural unit of the embarrassingly-parallel
algorithm; no aggregation is needed because per-point work is
uniform.

Real hardware fixes the native ratio, so the method exposes
an optional \texttt{gpu\_slowdown\_factor} that injects a
controlled post-kernel delay on the GPU worker, moving the
operating point along the ratio axis so the
ratio-versus-speedup relationship can be measured across
regimes the native workload set does not span. This is the
software analog of the GPU-throughput-sweep methodology
established by Mei and Chu~\cite{mei2017gpudvfs} (where the
native knob is DVFS); the sweep instrument is methodological,
not a scheduling mechanism. Native and achieved CPU/GPU
ratios are reported alongside each measurement.

The delay-injection knob is not a competing approach to
hardware-faithful GPU characterization. Validated
simulators~\cite{khairy2020accelsim}, kernel
sampling~\cite{avalos2021pka}, and model-driven throughput
predictions~\cite{gabbay2022dnnmemory} attribute runtime to
architectural causes inside the device; here the GPU is
treated as a black box and a controllable offset is added
to traverse the ratio axis without modeling kernel-level
contention. The purpose is to reach regimes the native
workload set does not span, not to explain why a given
throughput is observed.

\subsection{Adaptive Calibration with Cache Reuse
(RQ3)}\label{methods-calibration}

The hybrid scheduler (Section~\ref{methods-scheduling})
requires a per-machine calibration step to select its
hetero-config tuple: a policy choice $\in
\{\texttt{gpu\_only}, \texttt{gpu\_cpu\_hybrid}\}$ together
with GPU batch size, CPU worker count, and CPU batch size
when hybrid is selected. This calibration is the
unconditional upfront cost paid before any per-checkpoint
savings can be realized. A single-checkpoint user pays that
cost without reusing the selected tuple and loses; the cache
changes the unit of cost from per-checkpoint to per-session by
paying calibration once and reusing the selected tuple across
$N$ related checkpoints on the same machine. Loss-landscape
analysis is
rarely a single run: users compare related
checkpoints~\cite{xie2025losslens,losslens_repo}. A
\emph{related checkpoint} shares architecture, dataset, loss,
and grid spec but differs in learned weights (e.g.\ a
different seed).

\textbf{Composition with RQ1.} The GPU-side implementation
inside the hetero-config tuple is configured as RQ1's
per-(machine, workload) winner for the same workload; RQ1's
choices apply to the GPU worker only, since \texttt{vmap}-
and \texttt{torch.compile}-style transforms over the small
per-task CPU slices are out of scope. RQ1 is not in RQ3's
sweep — its winner is established by RQ1 and consumed here as
a configuration input. Composition with RQ1 happens at the
deployed-stack level, not the sweep level: the session
runtime assembles RQ1's winner with RQ3's calibrated
hetero-config tuple. The cache key (machine, workload)
matches RQ1's key so the composition is well-defined and the
amortization claim is keyed on the same identity.

\textbf{Calibration probe.} Calibration uses the smallest
square grid satisfying $r^2 \geq 4 \times (1 + p_{\max})$,
where $1 + p_{\max}$ is the total compute-unit count (one
GPU plus $p_{\max}$ CPU workers). The factor $4$ keeps the
per-unit task count large enough that scheduler granularity
is not the limiting factor of the probe. The probe sweeps
GPU batch size, CPU worker count, and CPU batch size with
patience-based
termination after \texttt{calibration\_retry}$=R$
consecutive non-improvements, inheriting OpenTuner's bounded
autotuning-search discipline~\cite{ansel2014opentuner} as the
departure from ATLAS-style exhaustive search. The selected
configuration is the lowest-runtime valid one, or
\texttt{gpu\_only} fallback if no hybrid configuration beats
the GPU-only baseline.

\textbf{Cache reuse.} For a related-checkpoint session, the
selected tuple is plausibly stable across checkpoints, so
the calibration cost paid by the first checkpoint is reused
for the rest. This instantiates the autotune-and-cache
pattern of ATLAS~\cite{whaley1998atlas} with the signature
lifted from op-level to workload-level: ATLAS keys on
(machine, BLAS op), RQ3 keys on (machine, workload =
architecture + dataset + loss + grid spec). The cached value
is the hetero-config tuple selection. Cache reuse is
invalid if the selected mode is not stable across the
session.

\textbf{Session comparison.} RQ3 is measured in the native
operating regime of the machine. The calibrated tuple is used
for every checkpoint in the composed session, and the
reference session runs the vanilla per-checkpoint path. The
reported quantities are the one-time calibration cost, the
cached per-checkpoint grid time, cumulative
$T_{\text{session}}$, and the break-even session size.

\subsection{Benchmarking, Surface Validation, and
Uncertainty}\label{methods-benchmarking}

\textbf{Timing primitive.} Elapsed time is measured at the
caller with \texttt{time.perf\_counter\_ns} as specified by
PEP 418~\cite{python_perf_counter_ns}: a portable monotonic
high-resolution clock that captures host-observed wall time
including Python orchestration, framework dispatch, backend
execution, parameter binding, dataloader iteration, result
collection, and synchronization. Asynchronous backend
completion is forced through the framework backend
abstraction before stopping the timer; no
CUDA-event-specific dependency is introduced. PyTorch Timer
and CUDA events are not root canon for this thesis.

\textbf{Runtime targets.} Two targets are named separately:
$T_{\text{grid}}$ is the wall time for one loss-grid after
inputs are ready (the RQ1 and RQ2 target), and
$T_{\text{session}}$ is the wall time for a multi-checkpoint
session, including any one-time calibration the measured arm
performs (the RQ3 target). Every speedup claim names the
target it uses. $T_{\text{grid}}$ is also reported in a
section decomposition $T_{\text{grid}} = T_{\text{perturb}} +
T_{\text{bind}} + T_{\text{eval}} + T_{\text{collect}} +
\text{residual}$ using the same wall-time primitive on each
section. The decomposition is explanatory only; it is not a
separate timing mechanism.

\textbf{Repeats and uncertainty.} Repetition and stopping
budget follow Kalibera and
Jones~\cite{kalibera2013rigorous}: repeated runs after
warm-up, variation estimates per arm, and a stopping budget
that bounds total experiment time. The speedup ratio is
reported with an effect-size confidence interval following
Kalibera and Jones~\cite{kalibera2020effect}. The verdict
rule is: $\text{CI}_{\text{low}} > 1.0 \Rightarrow$
supported speedup; $\text{CI}_{\text{high}} < 1.0 \Rightarrow$
supported regression; otherwise inconclusive. The same rule
applies across RQ1, RQ2, and — when full sessions are
repeated — RQ3; RQ3 is descriptive session evidence
otherwise.

\textbf{Trial order and pairing.} Baseline and candidate are
evaluated within the same repeat under identical workload,
seed, and grid; trial order rotates across repeats to avoid
fixed-order benchmarks. The pairing scheme is disclosed per
experiment.

\textbf{Surface-validation gate.} A candidate's timing claim
is conditional on producing the same loss surface as the
baseline on the same (workload $\times$ backend $\times$
checkpoint). Different implementations of the same nominal
float32 computation produce different numeric outputs in
practice (op reordering, kernel/fusion choice, hybrid
CPU/GPU split), so the gate is a relative-equivalence
criterion rather than bitwise equality. Three parts:
\begin{enumerate}
\item grid-point count and coordinates match the baseline
exactly;
\item paired NaNs at matching coordinates count as agreement;
\item finite losses satisfy $|a-b| \leq 1\mathrm{e}{-5}
\cdot \max(|a|,|b|)$ — $\text{rtol} = 1\mathrm{e}{-5}$
follows the \texttt{torch.allclose}
default~\cite{pytorch_allclose}; $\text{atol} = 0$, so
equivalence requires proportionally small drift rather than
an absolute floor.
\end{enumerate}
This gate is stricter than the visual-equivalence convention
used by canonical loss-landscape
implementations~\cite{li2018losslandscape,goldstein_loss_landscape_repo},
which validate landscape topology by inspection rather than
by per-point tolerance.
A surface that fails the gate disqualifies the timing claim
regardless of measured speedup.
